\subsection{Wave Optics} 

%\DTMnewdatestyle{dateupdated}{〈definition 〉

Answers to M\Alph{subsection} problems are on page \pageref{wave_optics_prob_answers}.  \textbf{\textcolor{red}{Updated \DTMnow.}}


%--------------------------------------------------------------------------------------------------------------
\begin{Exercise}
The distance between Earth and Mars varies between roughly 75~million kilometers and 375~million kilometers, depending on where they are in their respective orbits.  
Given the speed of light $c=3\times10^8$~m/s,  how long does it take a radio signal from Earth to reach one of the Mars rovers when the planets are at their closest point?
\end{Exercise}
\begin{Answer}
\end{Answer}



%--------------------------------------------------------------------------------------------------------------
\begin{Exercise}
The graph to the right shows a function $y(x)$ for which you know the exact value of both the function and its derivative $y'(x)$ at $x=10$:
$$
y(10)=37 \\
y'(10)=3
$$
(a) If $\Delta x=0.5$ in the picture, what is $\Delta y$?
(b) Estimate the value of $y(10.25)$.
(c) One of the following equations describes how to estimate $y(10.25)$ and the other three are total bogus. Which is the right one?
\begin{itemize}
\item $y(\Delta x)\approx \Delta y \cdot \frac{dy}{dx} \Delta x$
\item $y(\Delta x)\approx y(x)\Delta y + \frac{dy}{dx}y(x)$
\item $y(x+\Delta x)\approx y(x) + \frac{dy}{dx} \Delta x$
\item $y(x+\Delta y / \Delta x) \approx y(x) + \frac{dy}{dx}$
\end{itemize}


\end{Exercise}
\begin{Answer}
\end{Answer}



%--------------------------------------------------------------------------------------------------------------
\begin{Exercise}
The linear wave equation for a longitudinal sound wave along a thin rod is given by
$$
\frac{\partial^2D}{\partial x^2} = \frac{\rho}{E}\frac{\partial^2D}{\partial t^2},
$$
where $E$ is the \textit{Young's modulus} of the material---basically a measure of hardness---and $\rho$ is the mass density per volume of the material.  (The equation above is for the displacement $D(x,t)$ rather than the usual $y(x,t)$, because this is for a longitudinal wave; both the displacement and the velocity of the wave are in the $x$ direction, but you can't very well call it$ d^2 x/dx^2$, and $d^2 (Δx)/dx^2$ looks even worse.)  The hardest material known is diamond, which has a Young's modulus of $E=1.1\times10^12~\text{N/m}^2$ and a density of $\rho=3500~\text{kg/m}^3.  What is the speed of sound in a thin rod of diamond---assuming you could find a diamond rod long enough for anyone to care?  
\end{Exercise}
\begin{Answer}
\end{Answer}



%--------------------------------------------------------------------------------------------------------------
\begin{Exercise}
\end{Exercise}
\begin{Answer}
\end{Answer}



%--------------------------------------------------------------------------------------------------------------
\begin{Exercise}
\end{Exercise}
\begin{Answer}
\end{Answer}





\vspace{0.7in}

\textbf{Answers to Selected {\thesubsection} Problems:}
\label{wave_optics_prob_answers}
%\shipoutExercise
\shipoutAnswer

\cleardoublepage
